\renewcommand{\abstractname}{Abstrak}
\begin{abstract}

Ang komunikasyong panlipunan ay maaaring maging hamon sa mga "young adult," particular na sa kakayahang pragmatiko, o ang kakayahang 
mag-unawa ng implikadong kahulugan ng isang pahayag sa halip na unawaan ang literal na kahulugan lamang. Bagaman nagpapakita ng malakas 
na kakayahan ang mga artificial intelligence (AI) na chatbot sa pakikipag-usap, ang mga kasalukuyang kagamitan ay kadalasang nagkukulang 
sa kakayahang magbigay ng natural na simulasyon ng mga sitwasyong panlipunan. Higit pa rito ay madalas hindi rin naisasaalang-alang ang 
mga kultural na basehan at asal na nakapaloob sa mga sitwasyong ito para sa isang personal at ayon-sa-kontekstong pag-aaral ng 
komunikasyon. Dahil dito, binuo sa pag-aaral na ito ang Socius, isang "dual-agent conversational" AI na nakabatay sa "AI Partner-AI 
Mentor (APAM) Framework." Ginagamit ng sistema ang Google Gemini 3 API (Gemini 3-Flash) bilang mga AI agent para sa pakikipag-usap ng AI 
Partner at pagbibigay ng puna ng AI Mentor. Sa pag-aaral na ito ay nabuo rin ang isang "multi-pillar theoretical framework" na nagsasama 
ng Grice's Maxims, sikolohiyang pangkaunluran (Arnett at Havighurst) angkop sa mga young adult, at kulturang asal ng mga 
Pilipino (hal. Pakikisama at Hiya) upang makalikha ng mga kontekstong pang-"roleplay" kasama ang AI partner at kwantitatibong pagtatasa 
ng kakayahang pragmatiko at kwalitatibong puna mula sa AI Mentor. Mula sa "framework" na ito ay nabuo ang isang "synthesized scoring 
framework" na hango sa PragmEval Framework upang gabayan at markahan ang kakayahang pragmatiko ng mga gumagamit ng chatbot. Ang lahat ng 
bahaging ito ay ipinatupad bilang isang magkakaugnay na sistema upang magbigay ng isang "machine-learning driven na" personal na espasyo
para sa mga "emerging adult" upang magsanay at malinang ang kanilang kakayahang pragmatiko. Sinuri ang mga bahaging ito sa pamamagitan ng 
"functional testing", balidasyon ng eksperto, at "end-user testing" na binubuo ng 24 na "emerging adults." Ang mga pagsusuring 
ito ay nag-resulta sa sistema na gumagana ayon sa nilalayon sa pamamagitan ng "end-to-end testing," balidasyon ng eksperto sa ginamit na 
"theoretical framework," at ang pagturing ng mga end-user na ang sistema ay kawili-wili, kapaki-pakinabang, at epektibo. Ngunit, lumikha 
ito ng isang pedagohikal na paradoks sapagkat walang makabuluhang pagbabago sa kanilang kakayahang pragmatiko, batay sa 
"scoring framework." Inirerekomenda para sa mga susunod na pananaliksik ang pagpapalawig ng panahon ng ebalwasyon 
(longitudinal evaluation), pagpapatatag ng imprastraktura, sapilitang pakikipag-usap sa AI Mentor, at balidasyon ng mga eksperto sa 
kakayahan ng "scoring system" upang mas masusing masukat ang epektong pedagohikal ng Socius.

\begin{flushleft}
\begin{tabular}{lp{4.25in}}
\hspace{-0.5em}\textbf{Keywords:}\hspace{0.25em} & Komunikasyong panlipunan, Kakayahang pragmatiko, Conversational AI, Natural language processing, Teknolohiyang pantulong, Teknolohiyang edukasyonal, Grice’s Maxims, AI feedback systems
\end{tabular}
\end{flushleft}
\end{abstract}